\documentclass[10pt,a4paper]{amsart}
\usepackage[margin=19mm]{geometry}
\usepackage{amsmath,amssymb,mathtools}
\usepackage[scr=rsfs,cal=euler]{mathalfa}
\usepackage{xfrac}
\usepackage[unicode,hidelinks,bookmarks=false]{hyperref}
\newcommand{\ie}{i.e.\ }
\newcommand{\cf}{cf.\ }
\newcommand{\resp}{respectively\ }
\newcommand{\Subsection}{\subsection}
\providecommand{\nobreakdash}{\nobreak\hbox{-}}
\setlength{\emergencystretch}{2em}
\multlinegap0pt
\usepackage{bm}
\usepackage{bbm}
\usepackage{dsfont}
\usepackage{xspace}
\usepackage{cancel}
\usepackage{mathtools}
\usepackage{amsthm}
\makeatletter
\@ifundefined{theorem}{\newtheorem{theorem}{Theorem}}{}
\@ifundefined{lemma}{\newtheorem{lemma}[theorem]{\bf Lemma}}{}
\makeatother
\makeatletter
\expandafter\ifx\csname claimproof\endcsname\relax
\newenvironment{claimproof}[1][\proofname]{\proof[#1]\renewcommand{\qedsymbol}{$\blacksquare$}}{\endproof}
\fi
\expandafter\ifx\csname proof\endcsname\relax
\renewenvironment{proof}[1][\proofname] {\par\pushQED{\qed}\normalfont\topsep6\p@\@plus6\p@\relax\trivlist\item[\hskip\labelsep\bfseries#1\@addpunct{.}]\ignorespaces}{\popQED\endtrivlist\@endpefalse}
\fi
\makeatother
\title[Proof of the KAMAK tree conjecture]{Proof of the KAMAK tree conjecture}
\author{Micha Christoph, Raphael Steiner}
\date{Source: arXiv:2505.21367v1 (CC BY 4.0)}
\begin{document}
\maketitle

\noindent\emph{Source and licence.} Proof of the KAMAK tree conjecture. Version arXiv:2505.21367v1 (CC BY 4.0), distributed under the Creative Commons Attribution 4.0 International licence, as recorded in the arXiv licence metadata for this exact version and shown on the abstract page https://arxiv.org/abs/2505.21367v1. Full rights notice: licenses/arxiv-2505.21367-CC-BY-4.0.txt.\par\smallskip

\noindent\textit{Editorial scope.} Counted unit: the Lemma whose source label is \texttt{lem: lovasz trick} in the article (source0.tex lines 187--189, source label \texttt{lem: lovasz trick}), delivered together with the article's own complete original proof of it (source0.tex lines 190--215, one \texttt{proof} environment of 26 lines). No mathematics is written, repaired or re-derived: statement and proof are the article's own text line for line. Required context retained uncounted: lem: high degree (lines 160--162); lem: lovasz local lemma (lines 169--171); lem: get broom from degrees (lines 173--175). Every definition, display and label of those lines is retained unchanged. Omitted from this package and not redistributed: 1--159, 163--168, 172--172, 176--186, 216--413. Nothing inside a retained excerpt is omitted or altered: every retained body is delivered exactly as the article's lines are written, so no omission receipt of any kind accompanies this package. Citations: the retained excerpts carry no citation command at all, so no bibliography entry of the article is transcribed and no reference is redistributed. Reference adaptation: none was needed and none was made; every reference inside the delivered unit resolves inside it, so no unresolved reference or citation is produced. Printed slips kept literal: no printed word, symbol, hypothesis or number of the counted statement or of its proof was altered. Layout and class: the article's own class is \texttt{article} with packages including amsmath, amssymb, amsthm, graphicx, amscd, epic, eepic, url, xcolor, inputenc, comment, enumerate, epstopdf, enumitem; the wrapper is the lane's pinned \texttt{amsart} editorial wrapper at 10pt on A4, which restates the theorem-like environments the retained text needs and adds the article's own macro definitions that the retained text needs (none), with extra wrapper packages (bm, bbm, dsfont, xspace, cancel, mathtools). The delivered numbering is the unit's own. Assets: no retained span contains a figure environment, a graphics-inclusion command, a PDF-page inclusion, a table or a TikZ picture; no image file is copied into this package, the package declares no asset dependency and no image file is redistributed with it. Licence: version arXiv:2505.21367v1 of this article is distributed under the Creative Commons Attribution 4.0 International licence, as recorded in the arXiv licence metadata for this exact version and shown on the abstract page https://arxiv.org/abs/2505.21367v1. A full rights notice is delivered at licenses/arxiv-2505.21367-CC-BY-4.0.txt. Characters: every retained excerpt is pure ASCII, so the delivered engine, XeLaTeX with its default Latin Modern text font, reports no missing character.
\begin{lemma}\label{lem: high degree}
    Let $D$ be a $(k,d)$-broom digraph with root set $R$. Let $v\in V(D)$ and suppose that there is a directed path of length at most $k$ from $v$ to $R$. Then $d_D^+(v)= d$.
\end{lemma}

\begin{lemma}[Lov\'{a}sz Local Lemma]\label{lem: lovasz local lemma}
    Let $A_1,\ldots,A_m$ be a sequence of events such that each event occurs with probability at most $p$ and is independent of all the other events except at most $d$ of them. If $epd\leq 1$ then the event $\bigcap_{i=1}^{m}\overline{A_i}$ occurs with positive probability.
\end{lemma}

\begin{lemma}\label{lem: get broom from degrees}
    Let $k, d\in \mathbb{N}$ and let $T$ be an out-arborescence with root $r$ such that $d_T^+(r) \geq d$. Suppose that on every directed path in $T$ with vertex sequence $r=v_1,\ldots,v_h$ from $r$ to a leaf $v_h$, it holds that $d_T^+(v_i)\geq d$, for every index $i$ satisfying $\max\{1,h-k+1\}\leq i\leq h-1$. Then, $T$ contains a $(k-1,\lceil d/k\rceil)$-broom rooted at $r$ whose leaves are leaves of $T$.
\end{lemma}

\begin{lemma}\label{lem: lovasz trick}
    Let $d,k\in \mathbb{N}$ be such that $d\ge \max\{(4k)^6,5\cdot 10^{12}\}$ and let $D$ be a $(k,d)$-broom digraph with root set $R$. Then, there exists a $(k,\lfloor d^{1/(7k)}\rfloor)$-broom digraph $D'\subseteq D$ with root set $R'\subseteq R$ such that $d^{-}_D(r)\geq d^{1/10}$ for all $r\in R'$.
\end{lemma}

\begin{proof}
    Let $U:=\{u\in V(D)|d_D^+(u)=d\}$ and let $W:=\{w\in R|d_D^-(w)<d^{1/10}\}$. Note that by definition of a $(k,d)$-broom, we have $R\subseteq U$. For each arc $(u,v)\in A(D)$ with $u\in U$, let $X_{(u,v)}$ be a Bernoulli random variable that takes value $1$ with probability $d^{-2/3}$. For arcs $(u,v)$ with $u\notin U$, we set $X_{(u,v)}=1$. Note that by definition of a $(k,d)$-broom digraph, this does not include any arc ending in a vertex in $R$. In particular, for every arc $e$ ending in $W\subseteq R$ we have that $X_e=1$ with probability $d^{-2/3}$. Let $H\subseteq D$ denote the random subdigraph obtained by keeping only the arcs $e$ for which $X_e=1$.
    
    For each $u\in U$, let $A_u$ be the random event that $d^+_{H}(u)\leq d^{1/3}/2$. A simple application of Chernoff's inequality shows that $\mathbb{P}[A_u]\le \exp(-d^{1/3}/8)$. For each $w\in W$, let $B_w$ be the random event that $d^-_H(w)\geq 2$. By a union bound over all potential choices of two incoming edges, we obtain $\mathbb{P}[B_w]\le \binom{d^{1/10}}{2}d^{-4/3}\leq d^{-17/15}$. Furthermore, the events $A_u$ depend on at most $d$ events of the form $B_w$ and are independent of all other events. The events $B_w$ depend on at most $d^{1/10}\leq d$ events of the form $A_u$ and are independent of all other events. Using our assumption on $d$, one now easily checks that $epd\le 1$ for $p:=\max\{\exp(-d^{1/3}/8), d^{-17/15}\}$ and hence we can apply Lemma~\ref{lem: lovasz local lemma}, which implies that with positive probability none of the events $A_u, u\in U$ or $B_w, w \in W$ occur. 
    
    In the following, with a slight abuse of notation, we use $H$ to denote an \emph{instance} of the random subdigraph $H$ for which none of the mentioned events occur, that is, $d_H^+(u)>d^{1/3}/2$ for all $u\in U$ and $d_H^-(w)\le 1$ for all $w\in W$. Also note, that by definition of $H$, we have $d_H^+(x)=d_D^+(x)\ge 1$ for every $x\notin U$, and hence $\delta^+(H)\ge 1$.

    Let $R'$ denote the set of vertices with in-degree at least $2$ in $H$ and note that $R'$ is non-empty. Observe that $R'\subseteq R\setminus W$. For each vertex $r\in R'$, let us denote by $U_r$ be the set of vertices reachable with a directed path starting from $r$ in the digraph $H-(R'\setminus\{r\})$. Since all the vertices in $V(H)\setminus R'$ have in-degree at most $1$, it follows that the sets $U_r, r\in R'$ are pairwise disjoint. 
    
    For each $r\in R'$, let $D_r$ be the digraph obtained from $H[U_r]$ by removing all incoming arcs of~$r$. Observe that $D_r$ is an out-arborescence with root $r$, as $r$ has in-degree $0$ and all other vertices have in-degree $1$ and are reachable from $r$. Let $L_r\subseteq U_r$ be the set of vertices $u\in V(D_r)$ with $d_{D_r}^+(u)\leq d_H^+(u)/2$ and let $S_r\subseteq D_r$ denote the out-arborescence rooted at $r$ that is induced by all the vertices reachable from $r$ via directed paths in $D_r$ that do not use an arc starting in $L_r$. Note that $V(S_r)\cap L_r$ is exactly the set of leaves of $S_r$. 

    We claim that each leaf $l$ of $S_r$ must have (in $H$) at least one out-neighbor in $R$. Indeed, suppose not. Then $d_{H}^+(l)\ge 1$ since $\delta^+(H)\ge 1$, and at the same time $d_H^+(l)=|N_H^+(l)\setminus R'|= d_{D_r}^+(l)$. Therefore, $l\notin L_r$ which contradicts the assumption that $l$ is a leaf.
    
    Next, consider any vertex $u\in V(S_r)\subseteq U_r$. Then, we either have $u\in L_r$, in which case $u$ is a leaf of $S_r$ and satisfies $$|N_H^+(u)\cap (R'\setminus \{r\})|=d_H^+(u)-d_{H-(R'\setminus \{r\})}^+(u)=d_H^+(u)-d_{H[U_r]}^+(u)\ge d_H^+(u)-d_{D_r}^+(u)-1\ge \frac{d_H^+(u)}{2}-1,$$ or we have $u\notin L_r$, in which case, by definition of $S_r$ and $L_r$, we have $d_{S_r}^+(u)=d_{D_r}^+(u)> \frac{d_{H}^+(u)}{2}$.
    
    Recall that $D$ is a $(k,d)$-broom digraph and $S_r\subseteq D$. Let $v_1,\ldots,v_h$ be the vertex-sequence of any directed path $P$ from the root $r=v_1$ to a leaf $v_h$ in $S_r$. Since $v_h$ has an out-neighbor in $R$, it follows that, for $h-k+1\leq i\leq h$, $v_i$ has a directed path (in $D$) of length at most $k$ to a vertex in $R$. By Lemma~\ref{lem: high degree}, it follows that $d_D^+(v_i)=d$ and thus $v_i\in U$ by definition. Therefore, by our choice of $H$, we have that $d_{S_r}^+(v_i)>d_H^+(v_i)/2\ge d^{1/3}/4$ for $h-k+1\leq i\leq h-1$. Furthermore, since $v_h$ is a leaf of $S_r$ we have $v_h\in L_r$ and thus, by the above, $v_h$ has at least $d_H^+(v_h)/2-1\ge d^{1/3}/4-1$ out-neighbors in $R'\setminus \{r\}$. 
    
        Observe that by the freedom of choice of $P$ in this argument, we may also conclude that indeed every leaf of $S_r$ has at least $d^{1/3}/4-1$ out-neighbors in $R'\setminus \{r\}$.

    Let $d':=\lceil d^{1/3}/4\rceil$. 
    As we have shown that every vertex in $S_r$ at distance between $1$ and $k-1$ from a leaf has out-degree at least $d'$ in $S_r$, we can now apply Lemma~\ref{lem: get broom from degrees} to the out-arborescence $S_r$. This yields the existence of a $(k-1,\lceil d'/k\rceil)$-broom that is contained in $S_r$, rooted at $r$ and such that all its leaves are also leaves of $S_r$. One can easily check that our assumption on $d$ in the lemma implies that we have $d'/k\ge d^{1/(7k)}$, and hence we can, in fact, prune the above $(k-1,\lceil d'/k\rceil)$-broom to obtain a $(k-1,\lfloor d^{1/(7k)}\rfloor)$-broom $B_r$ contained in $S_r$ with root $r$ and all leaves are also leaves of $S_r$. 
    
    Note that it follows, by definition of a $(k-1,\lfloor d^{1/(7k)}\rfloor)$-broom, that $B_r$ has at most $(d^{1/(7k)})^k=d^{1/7}$ leaves. Since furthermore, by the above, every leaf of $B_r$ has at least $d^{1/3}/4-1$ out-neighbors in $R'\setminus \{r\}$, we may extend $B_r$ to a $(k,\lfloor d^{1/(7k)}\rfloor)$-broom $T_r$ by greedily choosing a subset of $\lfloor d^{1/(7k)}\rfloor$ out-neighbors in $R'\setminus \{r\}$ disjoint from all the other already chosen vertices for each leaf of $B_r$ (in some order). Note that this is indeed possible, since at any point in the process the current leaf has at least $d^{1/3}/4-1-d^{1/(7k)}\cdot (d^{1/7}-1)\ge d^{1/7k}$ remaining leafs to choose from, where the latter inequality follows by our assumptions on $d$ in the lemma.
    
    Note that the root and leaves of $T_r$ are in $R'$ while the internal vertices are all contained in $V(H)\setminus R'$. Let $D'\subseteq H$ be the union of all the brooms $(T_r)_{r\in R'}$ defined as above. Note that $D'$ is a $(k,\lfloor d^{1/(7k)}\rfloor)$-broom digraph with root set $R'\subseteq R$. Since $R'\subseteq R\setminus W$, we have $d_D^-(r)\ge d^{1/10}$ for every $r\in R'$, as desired. This concludes the proof of the lemma.
\end{proof}

\end{document}
